% dynamics/optimization_in_state_space.tex
\chapter{Optimization Produces Candidates}
\label{chap:optimization-in-state-space}

\section*{An objective must have an owner}

The preceding chapters can produce qualified model outputs and comparisons.
They do not make one quantity self-evidently desirable. Short transition
length, reduced balance residual, comfort, wear, cost and energy are distinct
objectives unless a responsible source explicitly aggregates them.

\begin{principle}[Qualified optimization principle]
\label{princ:state-space-optimization}
Optimization begins only after problem identity, purpose, decision variables,
fixed inputs, feasible set, model and objective have been declared.
\end{principle}

\begin{principle}[Candidate principle]
\label{princ:trajectory-shaping}
A solver returns a calculation candidate under that declaration, not an
Engineering Decision.
\end{principle}

\begin{figure}[htbp]
\centering
\begin{tikzpicture}[
  node distance=5mm and 6mm,
  box/.style={draw,rounded corners,align=center,minimum height=9mm,text width=27mm,font=\scriptsize},
  arrow/.style={->,thick}
]
\node[box] (problem) {problem identity\\purpose and freedoms};
\node[box,right=of problem] (model) {identified model \(M\)\\fixed inputs};
\node[box,right=of model] (objective) {declared \(J,\mathcal F\)\\units and source};
\node[box,below=of objective] (candidate) {calculation candidate\\values and residuals};
\node[box,left=of candidate] (evaluation) {qualified evaluation\\evidence and rules};
\node[box,left=of evaluation] (decision) {Engineering Decision\\authority and record};
\draw[arrow] (problem)--(model);
\draw[arrow] (model)--(objective);
\draw[arrow] (objective)--(candidate);
\draw[arrow] (candidate)--(evaluation);
\draw[arrow] (evaluation)--(decision);
\end{tikzpicture}

\caption{Dependency from an identified problem and model to a non-authoritative candidate, evaluation and possible Engineering Decision.}
\label{fig:model-candidate-decision}
\end{figure}

\section{Problem contract}

\begin{definition}[Runtime state trajectory]
\label{def:runtime-state-trajectory}
A runtime state trajectory is an identified model output
\(x_M(t)\) used by a declared optimization problem. It is not a universal AIM
state.
\end{definition}

\begin{definition}[Runtime state evolution]
\label{def:optimization-runtime-state-evolution}
Its evolution is generated by \(\mathcal D_M\) from fixed and variable inputs
under stated initial conditions and applicability.
\end{definition}

Only then is the familiar form meaningful:
\[
\min_{x\in\mathcal F}J(x).
\]
Here \(x\) denotes the declared decision vector, not an unspecified state.
The problem record shall identify:
\[
\begin{aligned}
\mathcal P=(\, &id,\ purpose,\ x,\ b,\ f,\ d,\ M,\ \mathcal F,\ J,\\
              &units,\ scaling,\ source,\ applicability,\ provenance\,),
\end{aligned}
\]
where \(b\) are bounds, \(f\) fixed inputs and \(d\) derived quantities.
Unresolved freedoms remain explicit rather than being filled silently.

\section{Degrees of freedom}

For a bounded continuity edit, a defensible example is
\[
x=(L,\alpha),\qquad
\mathcal F=\{x\mid c_{\mathrm{continuity}}(x)=0,
L_{\min}\le L\le L_{\max},
\alpha_{\min}\le\alpha\le\alpha_{\max}\}.
\]
Known endpoint states and the selected transition family are fixed;
\(L,\alpha\) are free within bounds; coefficients and terminal pose are
derived; continuity residuals are constrained. Cant, vehicle parameters and
an operational velocity profile remain unresolved unless deliberately added.

The current ufAIM/AXTRAN continuity solver is implementation evidence for
such bounded continuity solving. It is not evidence of general vehicle, cant,
comfort, network or automatic engineering-selection optimization.

\section{Objective and scaling}

\begin{definition}[Runtime-state functional]
\label{def:runtime-state-functional}
A runtime-state functional is a declared objective computed from an identified
model output, with units, scaling, source and purpose recorded.
\end{definition}

\begin{definition}[Operational trajectory shaping]
\label{def:optimization-trajectory-shaping}
Operational trajectory shaping is a candidate optimization formulation, not a
canonical AIM interpretation.
\end{definition}

The two objectives \(J_\kappa\) and \(J_r\) in
\cref{chap:vienna-reinterpretation} illustrate why scalar values cannot be
compared across purposes. A weighted aggregation
\[
J=w_\kappa J_\kappa/J_{\kappa,0}+w_rJ_r/J_{r,0}
\]
is meaningful only when the reference scales, weights and their responsible
source are recorded. The solver does not own them.

\section{Candidate contract}

A calculation candidate contains:
\[
\begin{aligned}
C=(\, &\mathcal P.id,\ M.id,\ x,\ b,\ \mathcal F,\ J,\ x^\star,\\
      &residuals,\ convergence,\ provenance,\ applicabilityStatus\,).
\end{aligned}
\]
The applicability status remains unevaluated on solver return, and the
candidate is explicitly non-authoritative.

\[
\boxed{\text{calculated}
\neq\text{applicable}
\neq\text{accepted}
\neq\text{authoritative}}
\]

Multiple candidates may be mathematically feasible, optimal under different
objectives, incomparable without purpose, invalid under later evidence, or
stale after an Alignment or realization change. One scalar score supplies no
universal engineering ordering.

\begin{definition}[Sparse runtime optimization]
\label{def:sparse-runtime-optimization}
Sparse runtime optimization is only a candidate computational technique that
exploits a stated sparse dependency structure; no general performance or
capability claim is made here.
\end{definition}

\begin{definition}[Operationally admissible trajectory]
\label{def:operationally-admissible-trajectory}
An operationally admissible trajectory is an evaluation result under
identified rules, purpose, evidence and applicability. Feasibility in
\(\mathcal F\) alone does not establish it.
\end{definition}

\section{Evaluation and Engineering Decision}

Evaluation asks whether a candidate and its evidence are applicable to a
purpose under identified rules, lifecycle and validity. An Engineering
Decision may then document selection or rejection together with authority.
The solver, AXTRAN implementation, App and Thesis do not hold that authority.

Thus the complete sequence is
\[
\begin{aligned}
S_{\mathrm{rail}}\to s(t)\to M\to x_M^{\mathrm{pred}}
&\to\text{comparison}\to J,\mathcal F\to C\\
&\to\text{evaluation}\to\text{Engineering Decision}.
\end{aligned}
\]
Geometry, model, objective, candidate and decision remain separately owned.
